\section*{Leitidee: Von lokalen �nderungen zu globalen Gr��en}\label{MULTI-leitidee}

Viele pyhsikalische, technische und wirtschaftliche Systeme h�ngen nicht nur von einer, sondern von mehreren Einflussgr��en ab - beispielsweise

\begin{itemize}
	\item Die Leistung eines Motors h�ngt von Drehzahl und Drehmoment ab.
	\item Die Dichte eines Materials kann von Ort und Temperatur abh�ngen.
	\item Die Energie eines Systems h�ngt von mehreren Zusatandvariablen ab.
\end{itemize}

Mathematisch beschreibt man solche Zusammenh�nge durch Funktionen mehrerer Ver�nderlicher

\[ f: \mathbb R^n \to \mathbb R, \vec x \mapsto f\left( \vec x\right)\,.\]

Die zentrale Idee der Analysis bleibt dabei dieselbe wie bei Funktionen einer Ver�nderlichen:

\begin{center}
	\textbf{Lokale �nderungen bestimmen das Verhalten des gesamten Systems.}
\end{center}

Die wichtigsten Werkzeuge daf�r sind

\begin{itemize}
	\item partielle Ableitungen,
	\item der Gradient,
	\item die lineare N�herung (Tangentialebene),
	\item Mehrfachintegrale.
\end{itemize}

Dieses Skript f�hrt in diese Grundideen der \textbf{multivariaten Analysis} ein.
